\chapter{Sprint 0 : Analyse et préparation}
\label{chap:sprint0}

\section*{Introduction}
Ce chapitre présente la phase de préparation du projet. Il regroupe l'identification des acteurs, la capture des besoins fonctionnels et non fonctionnels, le diagramme de contexte, le diagramme de cas d'utilisation global, le backlog du produit, la planification des releases et l'environnement technique utilisé.

\section{Analyse et spécification des besoins}

\subsection{Identification des acteurs}
La plateforme \textbf{LoomStay} est destinée à plusieurs types d'utilisateurs :
\begin{itemize}
    \item \textbf{Client} : accède à son espace chambre via QR code, effectue le check-in numérique, consulte les services de l'hôtel, envoie des messages et soumet des réclamations.
    \item \textbf{Réceptionniste} : gère les réservations, suit les check-ins, consulte les fiches voyageurs, communique avec les clients et gère les événements.
    \item \textbf{Administrateur} : gère les utilisateurs, les chambres, les informations de l'hôtel, les événements et consulte les journaux d'audit.
    \item \textbf{Responsable maintenance} : consulte les réclamations techniques, affecte les interventions aux employés, gère les shifts et suit l'avancement.
    \item \textbf{Gouvernante générale} : gère les réclamations housekeeping, affecte les tâches de ménage quotidien, gère les shifts et suit l'état des chambres.
    \item \textbf{Employé maintenance} : consulte les tâches assignées, scanne le QR code de la chambre, valide l'intervention et communique avec le responsable.
    \item \textbf{Employé housekeeping} : consulte les réclamations et les tâches de ménage quotidien, scanne le QR code et valide les tâches.
    \item \textbf{Système IA} : acteur technique secondaire appelé par le backend pour la détection de langue, la traduction automatique et la classification des réclamations.
\end{itemize}

\subsection{Diagramme de contexte du système}
\safefig{diagrams/contexte_global.png}{Diagramme de contexte de la plateforme LoomStay}{contexte_global}

\subsection{Capture des besoins fonctionnels}
\begin{table}[H]
\centering
\caption{Besoins fonctionnels de la plateforme LoomStay}
\small
\begin{tabularx}{\textwidth}{|c|X|p{3cm}|}
\hline
\textbf{ID} & \textbf{Besoin fonctionnel} & \textbf{Acteur(s)} \\ \hline
BF-01 & Authentification du personnel par email et mot de passe & Staff \\ \hline
BF-02 & Protection des routes selon le rôle utilisateur (RBAC) & Système \\ \hline
BF-03 & Gestion des chambres (CRUD, statut, étage, type) & Admin, Réception \\ \hline
BF-04 & Gestion des réservations (création, affectation chambre, check-in/out) & Réception \\ \hline
BF-05 & Génération de QR codes sécurisés (client et worker) & Système \\ \hline
BF-06 & Check-in numérique avec fiches voyageurs & Client \\ \hline
BF-07 & Suivi du taux de complétion du check-in & Réception \\ \hline
BF-08 & Accès à l'espace chambre via QR code avec validation du token & Client \\ \hline
BF-09 & Navigation dans l'espace chambre PWA installable & Client \\ \hline
BF-10 & Consultation des informations de l'hôtel & Client \\ \hline
BF-11 & Consultation des événements publiés avec images & Client \\ \hline
BF-12 & Convertisseur de devises vers le TND & Client \\ \hline
BF-13 & Messagerie bilingue en temps réel avec la réception & Client, Réception \\ \hline
BF-14 & Traduction automatique des messages (NLLB-200) & Système IA \\ \hline
BF-15 & Soumission de réclamations en texte libre & Client \\ \hline
BF-16 & Classification automatique des réclamations par IA & Système IA \\ \hline
BF-17 & Suivi, confirmation et réouverture des réclamations & Client \\ \hline
BF-18 & Affectation des réclamations aux employés & Responsable \\ \hline
BF-19 & Consultation des tâches dans l'application mobile Flutter & Employé \\ \hline
BF-20 & Scan QR entrée/sortie pour valider les interventions & Employé \\ \hline
BF-21 & Saisie du résultat et commentaire d'intervention & Employé \\ \hline
BF-22 & Messagerie interne liée aux réclamations & Employé, Manager \\ \hline
BF-23 & Gestion des tâches de ménage quotidien & Gouvernante \\ \hline
BF-24 & Gestion des shifts des employés & Responsable \\ \hline
BF-25 & Gestion des utilisateurs et profils employés & Admin \\ \hline
BF-26 & Gestion des événements et informations hôtelières & Admin, Réception \\ \hline
BF-27 & Consultation des journaux d'audit & Admin \\ \hline
\end{tabularx}
\label{tab:besoins_fonctionnels}
\end{table}
\FloatBarrier


\subsection{Capture des besoins non fonctionnels}
\begin{table}[H]
\centering
\caption{Besoins non fonctionnels de la plateforme LoomStay}
\small
\begin{tabularx}{\textwidth}{|c|p{2.8cm}|X|}
\hline
\textbf{ID} & \textbf{Catégorie} & \textbf{Description} \\ \hline
BNF-01 & Sécurité & Authentification par JWT, mots de passe hachés avec bcrypt (10~rounds), contrôle d'accès RBAC par middleware. \\ \hline
BNF-02 & Sécurité & Chiffrement AES-256-CBC des données sensibles (numéros de passeport) via le module \texttt{encryption.ts}. \\ \hline
BNF-03 & Sécurité & Tokens QR signés par JWT avec expiration automatique basée sur la date de checkout. \\ \hline
BNF-04 & Performance & Chargement paresseux (lazy loading) du modèle de traduction NLLB-200 pour optimiser la mémoire du service IA. \\ \hline
BNF-05 & Performance & Mise en cache des traductions fréquentes côté service IA pour réduire les appels au modèle. \\ \hline
BNF-06 & Performance & Communication temps réel via Socket.IO pour la messagerie, avec authentification par token sur chaque connexion. \\ \hline
BNF-07 & Utilisabilité & Interface responsive avec design sombre premium pour l'application Flutter. \\ \hline
BNF-08 & Utilisabilité & Espace chambre client installable en PWA avec manifest et service worker. \\ \hline
BNF-09 & Maintenabilité & Code source organisé en couches (Routes, Controllers, Services, Prisma) avec séparation des responsabilités. \\ \hline
BNF-10 & Maintenabilité & Schéma de base de données déclaratif avec Prisma ORM et migrations versionnées. \\ \hline
BNF-11 & Traçabilité & Journalisation des actions critiques via le modèle \texttt{AuditLog} avec métadonnées JSON. \\ \hline
BNF-12 & Fiabilité & Gestion du jour ouvrable hôtelier (business day) avec seuil à 06h00 pour les shifts et tâches quotidiennes. \\ \hline
\end{tabularx}
\label{tab:besoins_non_fonctionnels}
\end{table}
\FloatBarrier


\subsection{Diagramme de cas d'utilisation global}
\safefig{diagrams/usecase_global.png}{Diagramme de cas d'utilisation global simplifié}{usecase_global}

\section{Backlog du produit}
\begin{longtable}{|c|p{7.8cm}|p{2.4cm}|c|c|c|}
\caption{Backlog du produit LoomStay} \label{tab:backlog_produit} \\
\hline
\textbf{ID} & \textbf{User Story} & \textbf{Acteur} & \textbf{Prio.} & \textbf{Rel.} & \textbf{Spr.} \\
\hline
\endfirsthead
\hline
\textbf{ID} & \textbf{User Story} & \textbf{Acteur} & \textbf{Prio.} & \textbf{Rel.} & \textbf{Spr.} \\
\hline
\endhead
US-01 & Se connecter avec email et mot de passe & Staff & Must & R1 & S1 \\ \hline
US-02 & Protéger les accès selon les rôles & Système & Must & R1 & S1 \\ \hline
US-03 & Gérer les chambres de l'hôtel (CRUD, statut, étage, type) & Admin/Réception & Must & R1 & S2 \\ \hline
US-04 & Gérer les réservations (créer, modifier, affecter chambre) & Réception & Must & R1 & S2 \\ \hline
US-05 & Générer des QR codes sécurisés client et worker & Système & Must & R1 & S2 \\ \hline
US-06 & Effectuer le check-in numérique par numéro de réservation & Client & Must & R1 & S3 \\ \hline
US-07 & Remplir les fiches voyageurs (adultes et enfants) & Client & Must & R1 & S3 \\ \hline
US-08 & Suivre le taux de complétion du check-in & Réception & Must & R1 & S3 \\ \hline
US-09 & Accéder à l'espace chambre après validation du token QR & Client & Must & R2 & S4 \\ \hline
US-10 & Naviguer dans l'espace chambre PWA installable & Client & Must & R2 & S4 \\ \hline
US-11 & Consulter les informations de l'hôtel & Client & Should & R2 & S5 \\ \hline
US-12 & Consulter les événements publiés avec images & Client & Should & R2 & S5 \\ \hline
US-13 & Convertir des devises vers le TND & Client & Could & R2 & S5 \\ \hline
US-14 & Gérer les informations, événements et devises & Admin/Réception & Should & R2 & S5 \\ \hline
US-15 & Envoyer des messages à la réception & Client & Should & R2 & S6 \\ \hline
US-16 & Répondre aux messages clients & Réception & Should & R2 & S6 \\ \hline
US-17 & Traduire automatiquement les messages (NLLB-200) & Système IA & Must & R2 & S6 \\ \hline
US-18 & Envoyer une réclamation en texte libre & Client & Must & R3 & S7 \\ \hline
US-19 & Classifier la réclamation automatiquement par IA & Système IA & Must & R3 & S7 \\ \hline
US-20 & Suivre le statut d'une réclamation & Client & Must & R3 & S7 \\ \hline
US-21 & Confirmer ou rouvrir une réclamation & Client & Must & R3 & S7 \\ \hline
US-22 & Consulter et assigner les réclamations aux employés & Responsable & Must & R3 & S8 \\ \hline
US-23 & Transférer une réclamation entre services & Responsable & Should & R3 & S8 \\ \hline
US-24 & Créer et gérer les employés du département & Responsable & Must & R3 & S8 \\ \hline
US-25 & Consulter les tâches dans l'application mobile Flutter & Employé & Must & R3 & S8 \\ \hline
US-26 & Scanner le QR pour démarrer et terminer l'intervention & Employé & Must & R3 & S8 \\ \hline
US-27 & Rapporter le résultat de l'intervention (FIXED/NOT\_FIXED) & Employé & Must & R3 & S8 \\ \hline
US-28 & Échanger des messages internes liés aux réclamations & Employé/Manager & Should & R3 & S8 \\ \hline
US-29 & Gérer les tâches de ménage quotidien & Gouvernante & Must & R3 & S9 \\ \hline
US-30 & Gérer les shifts des employés & Responsable & Must & R3 & S9 \\ \hline
US-31 & Exécuter une tâche housekeeping/ménage avec scan QR & Employé HK & Must & R3 & S9 \\ \hline
US-32 & Consulter les journaux d'audit & Admin & Should & R3 & S9 \\ \hline
\end{longtable}


\section{Planification des releases}
\begin{table}[H]
\centering
\caption{Planification des releases}
\begin{tabularx}{\textwidth}{|c|X|c|}
\hline
\textbf{Release} & \textbf{Objectif} & \textbf{Sprints} \\ \hline
Release 1 & Socle système et gestion hôtelière : authentification, rôles, chambres, réservations, QR codes, check-in numérique & S1, S2, S3 \\ \hline
Release 2 & Parcours client et services numériques : espace chambre PWA, informations hôtel, événements, devise, messagerie bilingue & S4, S5, S6 \\ \hline
Release 3 & Réclamations, interventions et intelligence intégrée : réclamations IA, application Flutter, tâches housekeeping, shifts & S7, S8, S9 \\ \hline
\end{tabularx}
\label{tab:planification_releases}
\end{table}
\FloatBarrier


\section{Planification des sprints}
\begin{table}[H]
\centering
\caption{Répartition des sprints par release}
\small
\begin{tabularx}{\textwidth}{|c|c|X|c|}
\hline
\textbf{Rel.} & \textbf{Sprint} & \textbf{Contenu} & \textbf{Durée} \\ \hline
\multirow{3}{*}{R1} & Sprint 1 & Authentification, rôles et sécurité & 2 semaines \\ \cline{2-4}
                     & Sprint 2 & Gestion des chambres, réservations et QR codes & 2 semaines \\ \cline{2-4}
                     & Sprint 3 & Check-in numérique et fiches voyageurs & 2 semaines \\ \hline
\multirow{3}{*}{R2} & Sprint 4 & Espace client PWA et accès chambre & 2 semaines \\ \cline{2-4}
                     & Sprint 5 & Services hôteliers, événements et devise & 2 semaines \\ \cline{2-4}
                     & Sprint 6 & Messagerie client-réception et traduction & 2 semaines \\ \hline
\multirow{3}{*}{R3} & Sprint 7 & Réclamations clients et classification IA & 2 semaines \\ \cline{2-4}
                     & Sprint 8 & Affectation des réclamations et application Flutter & 2 semaines \\ \cline{2-4}
                     & Sprint 9 & Tâches housekeeping, ménage quotidien et shifts & 2 semaines \\ \hline
\end{tabularx}
\label{tab:sprints_par_release}
\end{table}
\FloatBarrier


\section{Environnement de travail}
\subsection{Technologies de développement}
\begin{table}[H]
\centering
\caption{Technologies de développement}
\small
\begin{tabularx}{\textwidth}{|p{2.5cm}|p{3cm}|X|}
\hline
\textbf{Catégorie} & \textbf{Technologie} & \textbf{Rôle} \\ \hline
Frontend web & React 18 + TypeScript & Construction de l'interface utilisateur web (staff et client) \\ \hline
Bundler & Vite & Serveur de développement et construction du frontend \\ \hline
Routing web & React Router v6 & Gestion des routes et de la navigation SPA \\ \hline
UI Components & shadcn/ui + Radix & Composants UI accessibles et personnalisables \\ \hline
Backend & Node.js + Express & Serveur API REST et logique métier \\ \hline
Langage backend & TypeScript & Typage statique et sécurité du code backend \\ \hline
ORM & Prisma & Mapping objet-relationnel et migrations PostgreSQL \\ \hline
Base de données & PostgreSQL & Système de gestion de base de données relationnel \\ \hline
Temps réel & Socket.IO & Communication bidirectionnelle (messagerie) \\ \hline
Service IA & FastAPI (Python) & API pour la traduction et la classification \\ \hline
Mobile & Flutter + Dart & Application mobile cross-platform pour les employés \\ \hline
Navigation mobile & GoRouter & Gestion des routes dans l'application Flutter \\ \hline
HTTP mobile & Dio & Client HTTP avec intercepteurs JWT \\ \hline
QR Scanner & mobile\_scanner & Scan des QR codes dans Flutter \\ \hline
\end{tabularx}
\label{tab:technologies_developpement}
\end{table}
\FloatBarrier


\subsection{Technologies complémentaires}
\begin{table}[H]
\centering
\caption{Technologies complémentaires}
\small
\begin{tabularx}{\textwidth}{|p{2.5cm}|p{3.5cm}|X|}
\hline
\textbf{Catégorie} & \textbf{Technologie} & \textbf{Rôle} \\ \hline
Authentification & JSON Web Token (JWT) & Gestion des sessions et des tokens d'accès \\ \hline
Hachage & bcrypt & Hachage sécurisé des mots de passe (10~rounds) \\ \hline
Chiffrement & AES-256-CBC & Chiffrement des données sensibles (passeports) \\ \hline
Traduction IA & NLLB-200-distilled-600M & Modèle pré-entraîné de traduction multilingue (Meta AI) \\ \hline
Classification IA & scikit-learn (TF-IDF + LR) & Pipeline de classification des réclamations \\ \hline
Détection langue & langdetect & Détection automatique de la langue des messages \\ \hline
Stockage & Supabase Storage & Stockage des images (événements) en démonstration \\ \hline
Versionnement & Git + GitHub & Gestion du code source et collaboration \\ \hline
Tests backend & Vitest & Framework de tests unitaires et d'intégration \\ \hline
Tests IA & Pytest & Tests unitaires du service IA FastAPI \\ \hline
Analyse mobile & flutter analyze & Analyse statique du code Dart \\ \hline
Stockage mobile & flutter\_secure\_storage & Stockage sécurisé du JWT sur l'appareil \\ \hline
Déploiement demo & Vercel / Render / HF Spaces & Hébergement de démonstration \\ \hline
\end{tabularx}
\label{tab:technologies_complementaires}
\end{table}
\FloatBarrier


\section{Architecture d'application}

\subsection{Architecture logique}
La solution repose sur une architecture orientée services composée de quatre modules principaux :
\begin{itemize}
    \item \textbf{Application web React/TypeScript} : interface pour le personnel (réception, administration, managers) et espace client chambre (PWA). Le frontend communique avec le backend via des API REST et Socket.IO pour le temps réel.
    \item \textbf{Backend Node.js/Express/TypeScript} : centralise la logique métier selon une architecture en couches (Routes $\rightarrow$ Controllers $\rightarrow$ Services $\rightarrow$ Prisma/PostgreSQL). Il expose les API REST, gère l'authentification JWT, le contrôle d'accès RBAC et communique avec le service IA.
    \item \textbf{Service IA FastAPI/Python} : fournit la détection de langue, la traduction via NLLB-200 et la classification des réclamations via TF-IDF / Régression Logistique. Il est appelé par le backend en tant que microservice.
    \item \textbf{Application mobile Flutter/Dart} : destinée aux employés terrain (maintenance et housekeeping). Elle permet de consulter les tâches, scanner les QR codes et valider les interventions.
\end{itemize}

\safefig{diagrams/architecture_generale.png}{Architecture générale de LoomStay}{architecture_generale}

\subsection{Architecture frontend orientée composants}
Le frontend React adopte une architecture orientée composants. Les pages sont organisées par rôle (\texttt{staff/} pour le personnel, \texttt{public/} pour les clients). Les routes sont protégées par le composant \texttt{AuthGuard.tsx} qui vérifie le rôle JWT. L'espace client chambre est protégé par \texttt{RoomLayout.tsx} qui valide le token QR. Les services API sont centralisés dans le dossier \texttt{api/}.

\subsection{Architecture backend en couches}
Le backend suit une architecture en couches clairement séparée :
\begin{enumerate}
    \item \textbf{Routes} : définissent les endpoints REST et appliquent les middlewares d'authentification et de rôle.
    \item \textbf{Controllers} : reçoivent les requêtes HTTP, valident les entrées et délèguent aux services.
    \item \textbf{Services} : contiennent la logique métier et interagissent avec Prisma ORM.
    \item \textbf{Prisma/PostgreSQL} : couche de persistance avec le schéma déclaratif et les migrations.
\end{enumerate}

\subsection{Architecture mobile}
L'application Flutter adopte une structure par fonctionnalités (\textit{feature-based}) :
\begin{itemize}
    \item \texttt{features/auth/} : authentification et stockage sécurisé du JWT.
    \item \texttt{features/tasks/} : liste des tâches, détails, scan QR, résultat d'intervention, messagerie.
    \item \texttt{features/home/} : tableau de bord de l'employé.
    \item \texttt{core/} : client API (Dio), routeur (GoRouter), constantes.
    \item \texttt{shared/} : widgets réutilisables.
\end{itemize}

\subsection{Architecture du service IA}
Le service IA FastAPI est structuré en trois services spécialisés :
\begin{itemize}
    \item \texttt{language\_service.py} : détection de langue avec \texttt{langdetect} et mapping vers les codes NLLB.
    \item \texttt{translation\_service.py} : traduction avec le modèle pré-entraîné \texttt{facebook/nllb-200-distilled-600M}, chargement paresseux et cache.
    \item \texttt{classifier\_service.py} : classification des réclamations avec un pipeline TF-IDF + Régression Logistique entraîné sur un jeu de données de 604~exemples en 6~catégories.
\end{itemize}

\subsection{Déploiement cible}
Le déploiement réel visé repose sur une infrastructure privée au sein de l'hôtel. Les plateformes Vercel, Render, Hugging Face Spaces et Supabase ont été utilisées uniquement pour la démonstration. Le chapitre~\ref{chap:deploiement} détaille l'architecture de déploiement cible.

\safefig{diagrams/deploiement_prive.png}{Architecture de déploiement cible sur serveur privé hôtelier}{deploiement_prive}

\subsection{Modèle de données global}
Le modèle de données global regroupe les 17~entités de la base PostgreSQL et leurs relations. Ce modèle est détaillé progressivement au fil des sprints.

\safefig{diagrams/modele_global.png}{Modèle de données global de LoomStay}{modele_global}

\section{Définition de Done}
Dans ce projet, un incrément de sprint est considéré comme terminé lorsque les fonctionnalités développées sont intégrées, testées, validées et documentées. Les tests, la construction de l'application et la validation sont des activités continues intégrées à chaque sprint, conformément aux principes Scrum. Le chapitre~\ref{chap:build_tests} constitue un chapitre de validation technique globale et non une release Scrum supplémentaire.

\section*{Conclusion}
Ce chapitre a permis de définir les besoins, les acteurs, le backlog, la planification et l'architecture générale de la solution. Le projet est organisé en trois releases principales, chacune contenant trois sprints fonctionnels.
